\documentclass[11pt]{article}

\usepackage[margin=1in]{geometry}
\usepackage{amsmath,amssymb,booktabs,tabularx}

\title{Pricing density normalization}
\author{}
\date{}

\begin{document}
\maketitle

\noindent\textbf{Status.} Draft\\
\textbf{Plan item.} Pricing density normalization\\
\textbf{Depends on.} \texttt{07\_terminal\_cumulant\_adapter.tex}\\
\textbf{Feeds.} Martingale tilt, risk-neutral characteristic function, terminal pricing measure

\section{Purpose}

Derive the normalized terminal density used to change from the physical
terminal law to the pricing terminal law.
The goal is to show that activity premia, signed-imbalance exposure risk price,
and the scalar return tilt define a valid likelihood ratio when the terminal
cumulant is finite.

\section{Setup}

Let
\[
    K_{t,T}(v,\gamma_+,\gamma_-,\rho_\Omega)
    =
    \log\mathbb E^P\!\left[
            \exp\!\left(
                vA_{t,T}^{\Omega}
                +vX_{t,T}^{\mathrm{core}}
                +\gamma_+N_{t,T}^+
                +\gamma_-N_{t,T}^-
                -\rho_\Omega B_{t,T}^\Omega
            \right)
\mid\mathcal F_t
\right]
\]
be the physical stochastic terminal cumulant produced by the adapter in
\texttt{07\_terminal\_cumulant\_adapter.tex}.
The deterministic terminal component \(d_{t,T}\) is not included in \(K_{t,T}\).
The market-pressure anchor-pressure movement \(A_{t,T}^{\Omega}\) is
stochastic and unweighted by mean-reversion core-shock weights.
\(B_{t,T}^\Omega\) is the signed-balance exposure integral whose drift-scaled
version is \(A_{t,T}^{\Omega}\).

Let \(\pi_+,\pi_->0\), and write
\[
    \gamma_+=\log\pi_+,
    \qquad
    \gamma_-=\log\pi_-.
\]

\section{Notation}

\begin{tabularx}{\textwidth}{lXX}
\toprule
Symbol & Meaning & Units or domain\\
\midrule
\(P\) & Physical measure & Base terminal law\\
\(Q^{\theta,\pi,\rho_\Omega}\) & Terminal pricing measure induced by the density & Equivalent terminal measure on the admissible domain\\
\(\theta\) & Scalar stochastic-return tilt & Real transform domain\\
\(\pi_\pm\) & Activity-premium count multipliers & Strictly positive\\
\(\gamma_\pm\) & Log activity premia & Real finite\\
\(M_T(\theta,\pi_+,\pi_-,\rho_\Omega)\) & Terminal density process value & Positive and unit conditional mean\\
\(\rho_\Omega\) & Signed-imbalance exposure risk price & Real finite\\
\(K_{t,T}\) & Physical stochastic cumulant & Finite on requested arguments\\
\bottomrule
\end{tabularx}

\section{Assumptions used}

\begin{tabularx}{\textwidth}{XXX}
\toprule
Assumption & Where it is used & Why it is needed\\
\midrule
The cumulant \(K_{t,T}(\theta,\gamma_+,\gamma_-,\rho_\Omega)\) is finite.
& Defining the normalizing constant.
& The likelihood ratio needs a finite conditional expectation.\\
The activity premia are strictly positive.
& Writing \(\gamma_\pm=\log\pi_\pm\) and preserving density positivity.
& Non-positive count multipliers cannot define a positive density.\\
\(\theta\) is real.
& Interpreting the exponential as a Radon--Nikodym density.
& Complex tilts are transform arguments, not probability measures.\\
\(\rho_\Omega\) is real and finite.
& Interpreting signed-imbalance exposure weighting as part of the density.
& Non-finite exposure risk prices cannot define a finite likelihood ratio.\\
The density is normalized conditionally on \(\mathcal F_t\).
& Proving unit conditional mean.
& A terminal pricing measure must preserve total probability.\\
\bottomrule
\end{tabularx}

\section{Result}

\noindent\textbf{Identity.} For every admissible real \(\theta\), strictly
positive \(\pi_\pm\), and finite \(\rho_\Omega\), define
\[
    M_T(\theta,\pi_+,\pi_-,\rho_\Omega)
    =
    \exp\!\left(
        \theta A_{t,T}^{\Omega}
        +\theta X_{t,T}^{\mathrm{core}}
        +\log(\pi_+)N_{t,T}^+
        +\log(\pi_-)N_{t,T}^-
        -\rho_\Omega B_{t,T}^\Omega
        -K_{t,T}(\theta,\log\pi_+,\log\pi_-,\rho_\Omega)
    \right).
\]
Then
\[
    \mathbb E^P\!\left[
        M_T(\theta,\pi_+,\pi_-,\rho_\Omega)\mid\mathcal F_t
    \right]
    =1.
\]

\section{Derivation}

By definition of \(K_{t,T}\),
\[
    \exp\!\left(K_{t,T}(\theta,\gamma_+,\gamma_-,\rho_\Omega)\right)
    =
    \mathbb E^P\!\left[
        \exp\!\left(
            \theta A_{t,T}^{\Omega}
            +\theta X_{t,T}^{\mathrm{core}}
            +\gamma_+N_{t,T}^+
            +\gamma_-N_{t,T}^-
            -\rho_\Omega B_{t,T}^\Omega
        \right)
        \mid\mathcal F_t
    \right].
\]
Set \(\gamma_\pm=\log\pi_\pm\).
The density is
\[
    M_T
    =
    \frac{
        \exp\!\left(
            \theta A_{t,T}^{\Omega}
            +\theta X_{t,T}^{\mathrm{core}}
            +\gamma_+N_{t,T}^+
            +\gamma_-N_{t,T}^-
            -\rho_\Omega B_{t,T}^\Omega
        \right)
    }{
        \exp\!\left(K_{t,T}(\theta,\gamma_+,\gamma_-,\rho_\Omega)\right)
    }.
\]
Taking the conditional physical expectation gives
\[
    \mathbb E^P[M_T\mid\mathcal F_t]
    =
    \frac{
        \mathbb E^P\!\left[
            \exp\!\left(
                \theta A_{t,T}^{\Omega}
                +\theta X_{t,T}^{\mathrm{core}}
                +\gamma_+N_{t,T}^+
                +\gamma_-N_{t,T}^-
                -\rho_\Omega B_{t,T}^\Omega
            \right)
            \mid\mathcal F_t
        \right]
    }{
        \exp\!\left(K_{t,T}(\theta,\gamma_+,\gamma_-,\rho_\Omega)\right)
    }
    =1.
\]

Since \(\theta\), \(\gamma_+\), \(\gamma_-\), and \(\rho_\Omega\) are real, and
since \(\pi_\pm>0\), the numerator is strictly positive.
Therefore, \(M_T\) is a positive normalized terminal density on the
admissible domain.

\section{Domain and validity}

The result is valid when
\(K_{t,T}(\theta,\log\pi_+,\log\pi_-,\rho_\Omega)\) is finite, the affine
schedule domain constraints hold for the corresponding adapted query,
\(\theta\) and \(\rho_\Omega\) are real, and \(\pi_\pm>0\).
The normalized density alone does not impose the martingale condition.
That condition determines which \(\theta\) is used for pricing.

\section{Consequences}

\begin{itemize}
    \item The denominator cumulant is the log normalizer of the pricing density.
    \item Activity premia are public pricing-kernel coordinates, while normalization is handled by subtracting \(K_{t,T}\).
    \item The density normalization guarantees \(\phi_Q(0)=1\) once the risk-neutral characteristic function is formed.
    \item A finite normalized density is necessary but not sufficient for no-arbitrage pricing.
\end{itemize}

\section{Checks}

\begin{tabularx}{\textwidth}{XXX}
\toprule
Check & Expected result & Evidence location\\
\midrule
Density denominator reuse
& The cumulant at \(\theta\) is evaluated and stored as the denominator for later frequency evaluations.
& \texttt{solve\_terminal\_affine\_martingale\_theta} in \texttt{crates/janiform-engine-pricing/src/spectral/pricer/terminal\_affine/core.rs}\\
Zero-frequency normalization
& The terminal-affine characteristic function at frequency zero is one.
& Test \texttt{terminal\_affine\_measure\_change\_normalizes\_characteristic\_and\_martingale}\\
Activity premium positivity
& Non-positive activity premia are rejected before pricing-kernel use.
& Test \texttt{pricing\_kernel\_rejects\_non\_positive\_premia}\\
\bottomrule
\end{tabularx}

\end{document}
